%! Dividing Polynomials — Unit 3, Lesson 3.3 — Lesson Plan
% UNITS 3+ PACKET-FIRST plan (LESSON_SHAPE.md §0a, 2026-11-14; engagement rules §1, 2026-09-27):
% a paper warm-up, then the guided notes — three rounds of definition → example → ladder → check
% and a Closing Practice — projected from a deck that mirrors them, then the DeltaMath homework
% started in class. 5 / 45 / 10. Teacher-facing; never handed out.
\documentclass[10pt]{article}
\usepackage{algebra2-article}
\usepackage{algebra2-boxes}
\usepackage{graphicx}   % -article does NOT load graphicx; needed for thumbnails/screenshots
\graphicspath{{images/}}
\newcommand{\TallMath}[1]{$\displaystyle #1\rule[-1.4em]{0pt}{3.2em}$}
\newcommand{\CourseName}{Algebra 2}
\newcommand{\MeetingLength}{60 minutes}
\newcommand{\UnitNumberName}{Unit 3: Polynomial Functions \quad}
\newcommand{\LessonNumberName}{Lesson 3.3: Dividing Polynomials}
\newcommand{\rung}[1]{{\color{goldacc}\ifcase#1\or$\star$\or$\star\star$\or$\star\star\star$\fi}}

\begin{document}
\begin{center}
    {\Large\bfseries \CourseName \\
    {\small\bfseries \UnitNumberName \LessonNumberName}}
\end{center}

\begin{tcolorbox}[colback=forestbg, colframe=forest, boxrule=0.9pt, arc=2mm,
    left=2mm, right=2mm, top=1.5mm, bottom=1.5mm]
    \textbf{Primary Objective:} Students divide a polynomial by a binomial (and by a quadratic)
    with long division, use synthetic division for a divisor $x - r$ with the correct $r$, write
    a zero placeholder for every missing power, and state and check every answer as
    $Q(x) + \frac{R}{D(x)}$ with $P = D\cdot Q + R$.
    \\[2pt]
    \textbf{Standards (2023 VA SOL):} \textbf{A2.EO.3c} --- divide polynomials (here binomial
    and quadratic divisors; monomial and factorable-trinomial divisors are left to the DeltaMath
    set if the teacher wants them); revisits \textbf{A2.EO.3a} (checking by multiplying).
    The Remainder and Factor Theorems are cut from the course and are \emph{not} named; a
    remainder of $0$ is read only as ``divides evenly.''
    \\[2pt]
    \textbf{Lesson model:} \emph{Packet first.} Students work in the printed packet the whole
    hour: a paper warm-up as they walk in; three rounds --- one merged skill each (the Unit 3
    exception) --- of a short \textbf{definition} and \textbf{example} ($\le$3 min), a
    \textbf{ladder} (\rung1\ \rung2\ \rung3), and a \textbf{Show me} or \textbf{Find the
    mistake}; a self-scored \textbf{Closing Practice}; and the homework started in class. The deck
    projects the same pages frame for frame. \textbf{Homework source:} DeltaMath online, set
    \emph{A2: Lesson 3.3: Dividing Polynomials}.
\end{tcolorbox}

\renewcommand{\arraystretch}{1.4}
\begin{skillbox}[Priority Ideas \& Skills]{goldbox}
\begin{tabularx}{\linewidth}{@{}X|X@{}}
\textbf{Priority Ideas \& Skills} & \textbf{Key Understandings (the ``why'')} \\[2pt]
\hline
\vspace{-6pt}
\begin{itemize}[leftmargin=1.1em, nosep, after=\vspace{2pt}]
  \item Divide with long division: divide the leading terms, multiply, \emph{subtract}, bring
        down; stop when the remainder's degree is below the divisor's.
  \item Write the result as $Q(x) + \frac{R}{D(x)}$ and check it with $P = D\cdot Q + R$.
  \item Divide synthetically by $x - r$, using $r$ (the sign of the constant flipped), and read
        a quotient one degree lower than the dividend.
  \item Write a $0$ for every missing power --- in the dividend, and in a long-division divisor
        --- before dividing.
\end{itemize}
&
\vspace{-6pt}
\begin{itemize}[leftmargin=1.1em, nosep, after=\vspace{2pt}]
  \item Division of polynomials is the same statement as division of numbers:
        $247 = 5\cdot49 + 2$ and $P = D\cdot Q + R$. The check is just Lesson 3.1's
        multiplication.
  \item Synthetic division is long division with the variables stripped out, so it only works
        when every column means one power of $x$ --- and only for $x - r$.
  \item \textbf{Target misconception (the crux, Round 3):} skipping the zero placeholder, so
        every coefficient after the gap lands in the wrong column and the quotient is wrong
        (often by a whole degree). \textbf{Flagged in Round 2:} using $-2$ instead of $2$ for
        $x - 2$.
\end{itemize}
\end{tabularx}
\end{skillbox}

\begin{skillbox}[{Vocabulary, Concepts, \& Theorems --- one definition each, worded as the notes print them}]{sky}
\small
\textbf{Polynomial long division} divides a polynomial $P(x)$ by a divisor $D(x)$ with the same
four steps as long division of numbers --- divide the leading terms, multiply, subtract, bring
down --- repeated until the remainder's degree is less than the divisor's; the answer is written
$\frac{P(x)}{D(x)} = Q(x) + \frac{R}{D(x)}$, and it always checks: $P(x) = D(x)\cdot Q(x) + R$.
\\[3pt]
\textbf{Synthetic division} is a shortcut for dividing by a divisor of the form $x - r$ that uses
only the coefficients: write $r$ --- the number that makes the divisor zero, so $x - 2$ uses $2$
and $x + 2 = x - (-2)$ uses $-2$ --- beside the dividend's coefficients, bring the first
coefficient down, then multiply by $r$ and add, column by column; the last number is the
remainder, and the numbers before it are the coefficients of the quotient, one degree lower.
\\[3pt]
A \textbf{zero placeholder} is a $0$ written for every power of $x$ that is missing: before you
divide, write the dividend (and a long-division divisor) in standard form with every power, so
$x^3 - 8$ becomes $x^3 + 0x^2 + 0x - 8$ with coefficients $1, 0, 0, -8$.
\end{skillbox}

\begin{fixedskillbox}[Lesson at a Glance --- \MeetingLength]{forestbg}
\renewcommand{\arraystretch}{1.25}
\begin{tabularx}{\linewidth}{@{}l c X X@{}}
\textbf{Phase} & \textbf{Min} & \textbf{Students} & \textbf{Teacher} \\
\hline
Warm-Up & 5 & Three quick wins on the warm-up page as they walk in; score themselves & Stand at the door; reveal at 3 min; land ``how do you \emph{find} the quotient?'' \\
Guided Notes \& Practice & 45 & Three rounds (13 min each): watch ($\le$3), \emph{climb the ladder} \rung1$\to$\rung2$\to$\rung3\ (8), then \emph{Show me} fingers or \emph{Find the mistake} (2); then the Closing Practice (6), scored out of 4 & $\le$3 min talk a round; circulate the whole ladder; reveal when most have \rung2; read the fingers; never rush Round 3's Find the mistake \\
Homework & 10 & Start the DeltaMath set alone & Launch the first problem aloud, then circulate \\
\end{tabularx}
\end{fixedskillbox}

\begin{skillbox}[Engagement --- how this hour keeps everyone solving]{goldbox}
\textbf{Nobody waits:} every round ends in a ladder --- \rung1\ everyone (a quick win),
\rung2\ everyone, \rung3\ early finishers (a non-monic divisor, a word problem, a quadratic
divisor). \textbf{Everyone commits, safely:} \emph{Show me} is 1--4 fingers on the count of
three, all at once, after Rounds 1 and 2. \textbf{Critique a stranger:} Round 3's \emph{Find the
mistake} puts the skipped placeholder on an anonymous Student A. \textbf{Their world:} study-hall
desks in equal rows (Round 2 \rung3) and the spring formal's dance floor (Closing bonus).
\textbf{Short clock:} every practice frame
carries its minutes. Praise the climb more than the answer.
\end{skillbox}

\begin{skillbox}[Warm-Up (5 min, on paper as they walk in)]{forestbg}
\textbf{1.} $247 \div 5 = 49$ R $2$, written $247 = 5\cdot49 + 2$ and $\frac{247}{5} = 49 +
\frac25$ --- the statement $P = D\cdot Q + R$ and the answer form, in numbers.
\textbf{2.} $\frac{6x^3}{2x} = 3x^2$, $\frac{-8x^2}{x} = -8x$, $\frac{x^4}{x} = x^3$;
$(x^3 + 5x^2) - (x^3 - 2x^2) = 7x^2$ and $(4x^2 - 3x) - (4x^2 + 8x) = -11x$ --- the
\emph{divide} and \emph{subtract} moves inside every step (the second subtraction is the sign
slip of Round 1). \textbf{3.} $(x - 2)(x^2 + 3x + 1) = x^3 + x^2 - 5x - 2$, so the quotient is
$x^2 + 3x + 1$ with remainder $0$. \textbf{Handoff:} ``Nobody will hand you the
multiplication. How do you \emph{find} the quotient?'' --- Divide. That is Round 1.
\end{skillbox}

\boxguard
\begin{skillbox}[The Lesson --- three rounds and Closing Practice (45 min)]{forestbg}
\begin{multicols}{2}
\textbf{Round 1 --- Long division (13 min).} Read the definition in one breath; point at the
numbers column on the deck ($247 = 5\cdot49 + 2$). Work $(x^3 + 4x^2 - 3x - 5) \div (x + 2) =
x^2 + 2x - 7 + \frac{9}{x+2}$; students fill the empty bracket on the page as you go. Say the
subtraction aloud --- ``the minus hits \emph{both} terms'' --- and close with the check.
\textbf{Ladder:} \rung1\ $x + 4$; \rung2\ $x^2 + x - 2 + \frac{3}{x-3}$; \rung3\ $3x^2 - 2x - 1 +
\frac{8}{2x+3}$ (non-monic divisor: $6x^3 \div 2x$). \textbf{Show me}
$(x^2 + 5x + 9) \div (x + 2)$: answer \textbf{2}, $x + 3 + \frac{3}{x+2}$. \textbf{1} dropped the
remainder; \textbf{3} subtracted only the leading term; \textbf{4} put $R$ over $x$.

\textbf{Round 2 --- Synthetic division (13 min).} Same example, faster: $r = -2$, bottom row
$1, 2, -7 \mid 9$ --- the same answer as Round 1, which is the point. Point at the ``which
$r$?'' table: \emph{flip the sign}. Say once that it works only for $x - r$ (\rung3\ of Round 1
could not use it). \textbf{Ladder:} \rung1\ $x + 5$; \rung2\ $2x^2 - x - 1 + \frac{2}{x-2}$;
\rung3\ $x^2 + x - 6$ desks a row ($r = -1$, remainder $0$: divides evenly). \textbf{Show me}
$(x^3 - 4x^2 + x + 6) \div (x - 2)$: answer \textbf{3}, $x^2 - 2x - 3$. \textbf{1} and
\textbf{4} used $-2$ for $x - 2$ (the Round 2 flag); \textbf{2} kept the dividend's degree.

\columnbreak
\textbf{Round 3 --- Missing terms: zero placeholders (13 min) --- the crux.} Show the
coefficient table: $x^3 - 8$ is $1, 0, 0, -8$. Work $(2x^4 - 3x^2 + 5) \div (x + 1)$: five
columns, $r = -1$, $2x^3 - 2x^2 - x + 1 + \frac{4}{x+1}$. \textbf{Ladder:} \rung1\ $x^2 + 2x + 4$;
\rung2\ $x^2 - 3x + 2 + \frac{3}{x+3}$; \rung3\ $2x + \frac{3x - 4}{x^2 + 1}$ (long division,
divisor $x^2 + 0x + 1$; stop when the degree drops below 2). \textbf{Find the mistake:} Student A
divided $(x^3 - 3x^2 + 4) \div (x - 2)$ with $1, -3, 4$ and got $x - 1 + \frac{2}{x-2}$. Fix: the
$0$ for the $x$-term; $1, -3, 0, 4$ gives $x^2 - x - 2$, remainder $0$ --- it divides evenly.
Do not rush this frame: let them find it before you reveal.

\textbf{Closing Practice (6 min).} Alone, scored out of 4. \textbf{1} $3x + 5 + \frac{5}{2x-3}$
(Round 1); \textbf{2} $x^2 + 4x - 1 + \frac{2}{x-2}$ (Round 2); \textbf{3} $(x^4 - 16) \div (x +
2) = x^3 - 2x^2 + 4x - 8$ (Round 3, the crux's twin: three zeros and $r = -2$); \textbf{4} yes,
$(x + 2)(2x - 3) + 2 = 2x^2 + x - 4$ (the check). \textbf{Bonus:} the formal's dance floor,
$2x^2 + 5x - 3$ ft. Say the \emph{Watch out} aloud: a $0$ for every missing power, and $x - 2$
uses $r = 2$. Collect the scores on fingers --- the exit read.
\end{multicols}
\end{skillbox}

\boxguard
\begin{skillbox}[Homework (DeltaMath) --- scored 1 to 5, started in class, due the first class after two study halls]{goldbox}
\textbf{Source:} DeltaMath, online --- set \textbf{A2: Lesson 3.3: Dividing Polynomials}; not in
the packet. \textbf{DeltaMath skills, in order:} \emph{Polynomial Long Division (Level 1)}; \emph{Polynomial Long Division (Level 2)}; \emph{Synthetic Division (Level 1)}; \emph{Synthetic Division (Level 2)}; \emph{Synthetic Division (Level 3)}. \textbf{Formative check:} \emph{Synthetic Division (Level 3)}. A missing-term synthetic division sorts
students into \emph{placeholder}, \emph{sign of $r$}, and \emph{secure}). Scored 1 to 5 from
DeltaMath; started in the last 10 minutes and due the first class after two study halls,
announced in class and on TurtleNet. \textbf{Preview:} Lesson 3.4 --- characteristics of
polynomial functions: degree, leading coefficient, end behavior, symmetry.
\end{skillbox}

\boxguard
\begin{skillbox}[Active Monitoring --- Watch For]{redbox}
\begin{itemize}[leftmargin=1.2em, nosep]
  \item \textbf{Subtracting only the first term} (Round 1 \rung2, Show me option 3): $-(x^3 -
        3x^2)$ written as $-x^3 - 3x^2$. Ask: ``what does the minus do to the second term?''
  \item \textbf{Wrong $r$} (Round 2 \rung1--\rung2, Show me options 1 and 4): $-3$ for $x - 3$.
        Have her say ``$x - 3$ is zero when $x = $\,\dots''
  \item \textbf{Skipped placeholder} (Round 3 \rung1--\rung2, Find the mistake, Closing 3): four
        columns for $x^3 - 8$ written as two. Point at the powers: ``where is your $x^2$?''
  \item \textbf{Quotient degree not lowered} (Show me option 2): reading $1, -2, -3$ as
        $x^3 - 2x^2 - 3x$.
  \item \textbf{Remainder over $x$, or dropped} (Round 1 Show me options 1 and 4).
  \item \textbf{Attitude:} a student parked at \rung1\ with nothing written --- crouch, point at
        the example on her page, ask for just the first step.
\end{itemize}
Cold-call prompts (after the fingers, never before): ``How do you know the division is
finished?'' \;``Why can't synthetic division do $2x + 3$?'' \;``How did Student A's answer lose a
degree?''
\end{skillbox}

% ── Teacher notes — packet order. The ONLY place teacher prose goes — never in a _key. ──
\boxguard[10]
\begin{teachernote}[Warm-Up]
On paper as they walk in; the deck's warm-up frame shows the same three items, revealed at three
minutes. Item 2's second subtraction ($-11x$) is the sign slip Round 1 will punish --- if many
miss it, say ``the minus reaches every term'' once before Round 1. Item 3 hands them a quotient
for free; the question it leaves (how to \emph{find} it) is Round 1's whole job.
\end{teachernote}

\boxguard[10]
\begin{teachernote}[Guided Notes \& Practice]
Three rounds of about 13 minutes and a 6-minute Closing Practice fill the 45. The examples'
brackets and synthetic frames print empty on the page --- students fill them as you work the
deck. Round 2's example is deliberately Round 1's problem again: the same answer by both methods
is what makes the shortcut believable. If behind, trim the \rung3s (they are for early finishers
anyway) and Round 1's Show me discussion --- never Round 3's Find the mistake, the crux. The
Remainder and Factor Theorems are cut: if a student notices the remainder equals $P(r)$, praise
the observation and move on.
\end{teachernote}

\boxguard[10]
\begin{teachernote}[Homework]
DeltaMath online, \emph{A2: Lesson 3.3: Dividing Polynomials}. Its skills, in order: \emph{Polynomial Long Division (Level 1)}; \emph{Polynomial Long Division (Level 2)}; \emph{Synthetic Division (Level 1)}; \emph{Synthetic Division (Level 2)}; \emph{Synthetic Division (Level 3)}. The last 10 minutes are for \emph{starting} it --- launch the first problem aloud, then
circulate toward anyone who scored under 3 on the Closing Practice. Due the first class after two
study halls; announce it in class and on TurtleNet.
\end{teachernote}

\boxguard[10]
\begin{teachernote}[Slides]
Seventeen frames that mirror the notes: warm-up, targets, then definition / example / ladder /
check for each round, the Closing Practice, and the homework frame. Every practice answer sits in
a click-to-reveal block --- reveal a ladder when most of the room has \rung2. The examples use
\verb|\pause| only to pace the steps. The 3-up slide handout is built as a stand-alone work
product, not part of the packet.
\end{teachernote}

\end{document}
